\clearpage
\SourceLeaf{356}
\pdfbookmark[2]{8. Gauss an Olbers, 29. Mai 1815}{parabolisch-008}
\begin{center}
[8.]\par
\medskip
\textsc{Gauss} an \textsc{Olbers}, Göttingen, 29. Mai 1815.
\end{center}
\noindent\hspace*{\parindent}.\ .\ .\ .\ .\ .\quad Ich wäre wohl geneigt, meine theoretischen Untersuchungen über die Berechnung der Kometenbahnen im Allgemeinen irgendwann in einem eigenen Werk als Supplement zu meiner Th. M. C. C. zu veröffentlichen. Es könnte vielleicht 6—8 Bogen umfassen, und ich würde dann noch eine neu eingerichtete Tafel für die parabolische Bewegung beifügen, deren Gebrauch noch etwas bequemer ist als der der \textsc{Barker}schen Tafel. Diese dürfte weitere 3 Bogen umfassen. Der Grund dafür ist schon gelegt, wobei mir Herr \textsc{Encke} geholfen hat.\quad .\ .\ .\ .\ .
\par\medskip
\begin{center}\rule{.18\textwidth}{.4pt}\end{center}
\medskip
\pdfbookmark[2]{9. Gauss an Bessel, 24. Juni 1815}{parabolisch-009}
\begin{center}
[9.]\par
\medskip
\textsc{Gauss} an \textsc{Bessel}, Göttingen, 24. Juni 1815.
\end{center}
\noindent\hspace*{\parindent}.\ .\ .\ .\ .\ .\quad Ich habe halb und halb die Absicht, demnächst ein Supplement zur Theoria motus corporum coelestium für die parabolischen Bahnen herauszugeben. Dazu habe ich mehrere theoretische Zusätze von einigem Gewicht. Auch denke ich eine neue Tafel für die parabolische Bewegung beizufügen, die noch bequemer eingerichtet ist als die \textsc{Barker}sche Tafel. Ihre Berechnung ist bereits so weit gediehen, dass ich für meinen eigenen Gebrauch schon jetzt keine andere mehr verwende.\quad .\ .\ .\ .\ .
\par\medskip
\begin{center}\rule{.22\textwidth}{.4pt}\end{center}
